% Thesis v3, Chuong 5 Discussion (cau truc Le duyet 12/11: 5.1 Interpretation, 5.2 Relation to Earlier Studies, 5.3 Threats, 5.4 Implications).
\chapter{Discussion}\label{ch:discussion}

This chapter interprets the empirical findings presented in Chapter~\ref{ch:results} in light of the evaluation-validity concerns identified in Chapter~\ref{ch:related-work}. The discussion examines what the combined results reveal about the evidential support for published PASS/FAIL decisions and how the observed patterns relate to earlier research. It also considers alternative explanations and threats to validity that may affect the strength and generalizability of these findings. On this basis, the chapter identifies implications for test-based evaluation practices, with recommendations grounded in the evidence and its remaining limitations.

\section{Interpretation of the Findings}\label{sec:d-interpretation}

Read together, the results of the three research questions point in one direction. The independent re-evaluation found substantially more evidence contradicting accepted decisions than evidence undermining rejections at the test stage, under the protocol applied to each branch. A retained divergence was found for \rqFPExeAllPct\% of the executable and \rqFPLibAllPct\% of the library acceptances that were fully re-evaluated, while \frc{rqFFExeAll} executable and \frc{rqFFLibAll} library rejections lack support. These proportions are not two measurements of one error rate. The accepted branch searches generated inputs for a single divergence, and the rejected branch replays the vectors that failed, so each proportion is tied to the way its evidence was gathered.

One reading that fits this pattern follows from what each kind of decision rests on. A rejection rests on a failing vector, which is an observed difference from an expected result, and the re-evaluation only has to confirm that the difference also holds against the C program. An acceptance rests on the absence of a failure on the available vectors, which does not establish behavioral agreement on untested inputs. The source of the inputs is consistent with this reading: every translation in $D$ has a retained divergence on a generated input, and only \rqDivTransLayerSExeAllNum{} executable translations have one on an original vector. The study does not test this reading against alternatives. It remains an interpretation of the pattern, and the difference between the two protocols may account for part of it.

A property of the corpus bears directly on the result of RQ2. TRACTOR checks in its continuous integration that every test vector matches the behavior of the C program~\cite{okhravi2026tractorbenchmark}. A vector on which a translation fails under the runner therefore already indicates a difference from the C program as the runner compares it, and an unsupported rejection requires either a comparison of the harness that differs from that of the runner or a behavior that changes between the two executions. Section~\ref{sec:r-rq2} is consistent with the first source: the one unsupported rejection depends on the rule for floating-point fields, and \rqRobFFNoStderrExeAllNum{} executable rejections become unsupported when standard error is ignored. The result of RQ2 therefore describes a corpus whose vectors are checked against the C programs, and it cannot be assumed to hold for benchmarks whose expectations are written independently of the reference program.

For the published success rates, the consequence is a decrease whose size is only partly determined. At the level of all systems, both intervals lie below the published rate. For executables the rate decreases by between \rqSRDecrMinExeAllPct{} and \rqSRDecrMaxExeAllPct{} percentage points, and the interval is \widthE{} points wide, so the smallest decrease is larger than the uncertainty that the open decisions leave. For libraries the rate decreases by between \rqSRDecrMinLibAllPct{} and \rqSRDecrMaxLibAllPct{} points, and the interval is \rqSRWidthTightLibAllPct{} points wide. Here the open decisions outweigh the smallest decrease: the data show that the pooled rate is lower after reclassification, and they leave the size of the decrease largely open. At the level of single systems the decrease is less uniform. The entire interval lies below the published rate for all ten systems on executables and for eight systems on libraries. Of the other two library systems, Gemini 3.1 Pro has an interval that exceeds the published rate because of decisions without a result, and SmartC2Rust has a published rate of zero.

The difference between the two kinds of programs needs the most careful reading. The proportion of RQ1 is \rqFPLibAllPct\% for libraries against \rqFPExeAllPct\% for executables, and the lower value does not show that library translations preserve behavior better. Two limits of the measurement weigh more heavily on libraries. The mutant-based validation detected \sMutLibK{} of \sMutLibN{} non-equivalent library mutants, compared with \sMutExeK{} of \sMutExeN{} executable mutants, which indicates a lower detection sensitivity for libraries under these controls, although the mutants come from only \sMutLibNCases{} library and \sMutExeNCases{} executable pilot cases. The sanitizers report undefined behavior on the C side for \rqUbDivLibAllPct\% of the library pairs of a translation and an input that show a divergence and for \rqUbDivExeAllPct\% of the executable ones. These percentages are counted over such pairs, and they show that a larger part of the observed differences is set aside for libraries. Both limits reduce the number of retained divergences that can be found for libraries. The study cannot separate their contributions from a real difference between the two kinds of programs, so the comparison between the two proportions remains open.

\section{Relation to Earlier Studies}\label{sec:d-relation}

On the side of acceptances, the result agrees in direction with what earlier studies observed when accepted outputs were tested further. In C-to-Rust translation, Li et al.\ ran fuzzing against the C programs on translations written by study participants and documented behavioral discrepancies in them \cite{li2025userstudy}. In code generation, EvalPlus showed that augmented tests lower the reported pass rates \cite{liu2023evalplus}, and in issue solving, PatchDiff and SWE-ABS found that a share of the patches accepted by SWE-bench Verified differ from the reference patch or fail under strengthened tests \cite{wang2026patchdiff,yu2026sweabs}. The proportion of RQ1 belongs to this family of findings: decisions of acceptance reached on fixed test suites have repeatedly been contradicted once the accepted outputs were exercised on further inputs. The agreement concerns the direction only. The studies differ in the artifacts they examine, in the reference against which behavior is compared and in the effort spent on generating tests, so their percentages cannot be set beside the \rqFPExeAllPct\% and \rqFPLibAllPct\% of this thesis. One distinction made by PatchDiff also applies here: a differential test establishes a behavioral difference, and a verdict of incorrectness needs further analysis, which that study carried out by hand and this thesis did not.

On the side of rejections, earlier evidence is sparse and points to the evaluation as a possible source of failures. Rabbi et al.\ inspected a subset of failing translations by hand and traced some of them to the evaluation setup, such as build commands, time limits and converted tests \cite{rabbi2026beyond}, and issue reports on CRUST-Bench describe test expectations that conflict with the behavior of the reference C program, so that a translation behaving like the C program fails \cite{crustbench-issues}. Against this background, the result of RQ2 is that almost none of the rejections that reached an outcome lack support. The two observations concern different populations and do not contradict each other. The failures reported earlier include problems of build configuration, while RQ2 examines rejections at the test stage only and leaves rejections at the build stage untouched. The result is also bounded by its own coverage: it holds for the rejections that reached an outcome on one corpus, and it moves from \frc{rqFFExeAll} to \frc{rqRobFFNoStderrExeAll} for executables when standard error is left out of the comparison.

The thesis differs from these studies in the object about which it draws conclusions. The earlier work establishes the value of independent testing or revises the result of one benchmark on the side of acceptances. This thesis takes the published PASS/FAIL decisions of several C-to-Rust systems as its object, adjudicates acceptances and rejections under explicit rules for each, and carries the decisions without a result into the reported rate as an interval. PatchDiff and SWE-ABS report point estimates of the revised rates and do not carry unexamined decisions into them \cite{wang2026patchdiff,yu2026sweabs}; the interval of RQ3 makes visible how much of the revised rate depends on decisions without a result. The price of this design is breadth. The thesis covers one corpus and ten systems, and its classification of divergences follows fixed rules where PatchDiff used manual analysis, so it supports statements about the evidence behind decisions and leaves the correctness of single translations open.

\section{Threats to Validity}\label{sec:d-threats}
\IfFileExists{generated/robust-numbers.tex}{% SINH TU DONG boi writing/scripts/gen-robust-macros.py. KHONG SUA TAY.
% Repo commit: a8b0b65a41f772c806ef4cc799c5ca7dd8831df7. Phan tich post-freeze tren S0 (exp 30 L3, exp 30 ctype post-hoc, exp 31, exp 32, exp 34). L3, B1, N1: kiem tren mau (R9, RQ2); B3: nhan dien giai, Le chap nhan khong kiem tay 2026-11-12.
\providecommand{\rqSet}[2]{\expandafter\def\csname #1\endcsname{#2}}
\rqSet{rqLThreeExeAllNum}{188}
\rqSet{rqLThreeExeAllDen}{504}
\rqSet{rqLThreeExeAllPct}{37.3}
\rqSet{rqLThreeNBLoExeAllNum}{163}
\rqSet{rqLThreeNBLoExeAllDen}{504}
\rqSet{rqLThreeNBLoExeAllPct}{32.3}
\rqSet{rqLThreeNBHiExeAllNum}{186}
\rqSet{rqLThreeNBHiExeAllDen}{504}
\rqSet{rqLThreeNBHiExeAllPct}{36.9}
\rqSet{rqBOneExeAllNum}{173}
\rqSet{rqBOneExeAllDen}{476}
\rqSet{rqBOneExeAllPct}{36.3}
\rqSet{rqBOneLeaveExeAllNum}{3}
\rqSet{rqBOneUncheckedExeAllNum}{0}
\rqSet{rqBOneInputsExeAllNum}{247}
\rqSet{rqBOneInputsExeAllDen}{4961}
\rqSet{rqBOneInputsExeAllPct}{5.0}
\rqSet{rqBThreeDYesExeAllNum}{90}
\rqSet{rqBThreeCasesYesExeAllNum}{17}
\rqSet{rqBThreeDNoExeAllNum}{85}
\rqSet{rqBThreeCasesNoExeAllNum}{19}
\rqSet{rqBThreeDUnclearExeAllNum}{1}
\rqSet{rqBThreeCasesUnclearExeAllNum}{1}
\rqSet{rqBThreeNotUbExeAllNum}{85}
\rqSet{rqBThreeNotUbExeAllDen}{476}
\rqSet{rqBThreeNotUbExeAllPct}{17.9}
\rqSet{rqRNineExeN}{60}
\rqSet{rqRNineExeUnclear}{0}
\rqSet{rqRNineFalseExeNum}{0}
\rqSet{rqRNineFalseExeDen}{60}
\rqSet{rqRNineFalseExePct}{0.0}
\rqSet{rqRNineFalseExeLo}{0.0}
\rqSet{rqRNineFalseExeHi}{6.0}
\rqSet{rqRNineBnuExeNum}{0}
\rqSet{rqRNineBnuExeDen}{60}
\rqSet{rqRNineBnuExePct}{0.0}
\rqSet{rqRNineBnuExeLo}{0.0}
\rqSet{rqRNineBnuExeHi}{6.0}
\rqSet{rqRNineWorstExeNum}{0}
\rqSet{rqRNineWorstExeDen}{60}
\rqSet{rqRNineWorstExePct}{0.0}
\rqSet{rqREighteenExeN}{9}
\rqSet{rqREighteenFalseExeNum}{0}
\rqSet{rqLThreeLibAllNum}{66}
\rqSet{rqLThreeLibAllDen}{1156}
\rqSet{rqLThreeLibAllPct}{5.7}
\rqSet{rqLThreeNBLoLibAllNum}{57}
\rqSet{rqLThreeNBLoLibAllDen}{1156}
\rqSet{rqLThreeNBLoLibAllPct}{4.9}
\rqSet{rqLThreeNBHiLibAllNum}{60}
\rqSet{rqLThreeNBHiLibAllDen}{1156}
\rqSet{rqLThreeNBHiLibAllPct}{5.2}
\rqSet{rqBOneLibAllNum}{65}
\rqSet{rqBOneLibAllDen}{1088}
\rqSet{rqBOneLibAllPct}{6.0}
\rqSet{rqBOneLeaveLibAllNum}{0}
\rqSet{rqBOneUncheckedLibAllNum}{0}
\rqSet{rqBOneInputsLibAllNum}{0}
\rqSet{rqBOneInputsLibAllDen}{1071}
\rqSet{rqBOneInputsLibAllPct}{0.0}
\rqSet{rqBThreeDYesLibAllNum}{21}
\rqSet{rqBThreeCasesYesLibAllNum}{7}
\rqSet{rqBThreeDNoLibAllNum}{44}
\rqSet{rqBThreeCasesNoLibAllNum}{21}
\rqSet{rqBThreeDUnclearLibAllNum}{0}
\rqSet{rqBThreeCasesUnclearLibAllNum}{0}
\rqSet{rqBThreeNotUbLibAllNum}{44}
\rqSet{rqBThreeNotUbLibAllDen}{1088}
\rqSet{rqBThreeNotUbLibAllPct}{4.0}
\rqSet{rqRNineLibN}{62}
\rqSet{rqRNineLibUnclear}{8}
\rqSet{rqRNineFalseLibNum}{0}
\rqSet{rqRNineFalseLibDen}{54}
\rqSet{rqRNineFalseLibPct}{0.0}
\rqSet{rqRNineFalseLibLo}{0.0}
\rqSet{rqRNineFalseLibHi}{6.6}
\rqSet{rqRNineBnuLibNum}{0}
\rqSet{rqRNineBnuLibDen}{54}
\rqSet{rqRNineBnuLibPct}{0.0}
\rqSet{rqRNineBnuLibLo}{0.0}
\rqSet{rqRNineBnuLibHi}{6.6}
\rqSet{rqRNineWorstLibNum}{8}
\rqSet{rqRNineWorstLibDen}{62}
\rqSet{rqRNineWorstLibPct}{12.9}
\rqSet{rqREighteenLibN}{9}
\rqSet{rqREighteenFalseLibNum}{0}
\rqSet{rqRTwoCheckExeN}{30}
\rqSet{rqRTwoCheckExeUnclear}{0}
\rqSet{rqRTwoCheckWrongExeNum}{0}
\rqSet{rqRTwoCheckWrongExeDen}{30}
\rqSet{rqRTwoCheckWrongExePct}{0.0}
\rqSet{rqRTwoCheckWrongExeLo}{0.0}
\rqSet{rqRTwoCheckWrongExeHi}{11.3}
\rqSet{rqRTwoCheckLibN}{31}
\rqSet{rqRTwoCheckLibUnclear}{0}
\rqSet{rqRTwoCheckWrongLibNum}{0}
\rqSet{rqRTwoCheckWrongLibDen}{31}
\rqSet{rqRTwoCheckWrongLibPct}{0.0}
\rqSet{rqRTwoCheckWrongLibLo}{0.0}
\rqSet{rqRTwoCheckWrongLibHi}{11.0}
\rqSet{rqNOneLibAllOnlyNum}{7}
\rqSet{rqNOneLibAllNum}{58}
\rqSet{rqNOneLibAllDen}{1088}
\rqSet{rqNOneLibAllPct}{5.3}
\rqSet{rqCtypeOutExeAllNum}{173}
\rqSet{rqCtypeOutExeAllDen}{469}
\rqSet{rqCtypeOutExeAllPct}{36.9}
\rqSet{rqCtypeOutNBLoExeAllNum}{155}
\rqSet{rqCtypeOutNBLoExeAllDen}{469}
\rqSet{rqCtypeOutNBLoExeAllPct}{33.0}
\rqSet{rqCtypeOutNBHiExeAllNum}{171}
\rqSet{rqCtypeOutNBHiExeAllDen}{469}
\rqSet{rqCtypeOutNBHiExeAllPct}{36.5}
\rqSet{rqCtypeCaseDExeNum}{3}
\rqSet{rqCtypeCasePtestExeNum}{7}
\rqSet{rqCtypeOutLibAllNum}{65}
\rqSet{rqCtypeOutLibAllDen}{1088}
\rqSet{rqCtypeOutLibAllPct}{6.0}
\rqSet{rqCtypeOutNBLoLibAllNum}{56}
\rqSet{rqCtypeOutNBLoLibAllDen}{1088}
\rqSet{rqCtypeOutNBLoLibAllPct}{5.1}
\rqSet{rqCtypeOutNBHiLibAllNum}{59}
\rqSet{rqCtypeOutNBHiLibAllDen}{1088}
\rqSet{rqCtypeOutNBHiLibAllPct}{5.4}
\rqSet{rqCtypeCaseDLibNum}{0}
\rqSet{rqCtypeCasePtestLibNum}{0}
}{}
\IfFileExists{generated/review-numbers.tex}{% SINH TU DONG boi writing/scripts/gen-review-macros.py. KHONG SUA TAY.
% Repo commit: cdab2263f99dd8f088c1181b64242790226feb6f. Nguon: experiments/40-review-counts/results/counts.json, experiments/39-prereg-backfill/results/v3c-{exe,lib}.json. Post hoc, Le chap nhan khong kiem tay 2026-11-12; khong thay so chinh.
\providecommand{\rqSet}[2]{\expandafter\def\csname #1\endcsname{#2}}
\rqSet{rqAgreeCUnstableVecExeNum}{0}
\rqSet{rqAgreeCUnstableVecExeDen}{0}
\rqSet{rqAgreeCUnstableVecExePct}{n/a}
\rqSet{rqAgreeCUnstableFTransExeNum}{0}
\rqSet{rqAgreeCUnstableFTransExeDen}{0}
\rqSet{rqResOnlyDivExeNum}{0}
\rqSet{rqResOnlyDivExeDen}{6655127}
\rqSet{rqResOnlyDivExePct}{0.0}
\rqSet{rqResMemOneSideDivExeNum}{0}
\rqSet{rqResXfszDivExeNum}{0}
\rqSet{rqResOnlyTransExeNum}{0}
\rqSet{rqResOnlyTransExeDen}{176}
\rqSet{rqResOnlyTransExePct}{0.0}
\rqSet{rqResAnyTransExeNum}{0}
\rqSet{rqVThreeCExeUnitsNum}{189}
\rqSet{rqVThreeCExeRetainedNum}{6827678}
\rqSet{rqVThreeCExeViolNum}{0}
\rqSet{rqAgreeCUnstableVecLibNum}{0}
\rqSet{rqAgreeCUnstableVecLibDen}{2}
\rqSet{rqAgreeCUnstableVecLibPct}{0.0}
\rqSet{rqAgreeCUnstableFTransLibNum}{0}
\rqSet{rqAgreeCUnstableFTransLibDen}{1}
\rqSet{rqResOnlyDivLibNum}{22}
\rqSet{rqResOnlyDivLibDen}{18194}
\rqSet{rqResOnlyDivLibPct}{0.1}
\rqSet{rqResMemOneSideDivLibNum}{22}
\rqSet{rqResXfszDivLibNum}{0}
\rqSet{rqResOnlyTransLibNum}{0}
\rqSet{rqResOnlyTransLibDen}{65}
\rqSet{rqResOnlyTransLibPct}{0.0}
\rqSet{rqResAnyTransLibNum}{3}
\rqSet{rqVThreeCLibUnitsNum}{589}
\rqSet{rqVThreeCLibRetainedNum}{18198}
\rqSet{rqVThreeCLibViolNum}{0}
\rqSet{rqVThreeCLibNoDivFileNum}{2}
}{}

The interpretation of Section~\ref{sec:d-interpretation} rests on five assumptions, and each of them can be questioned. This section takes the questions in turn. For each, it states which conclusion is at stake, what shows that the risk is real, what the study did about it, and what remains uncertain.

\paragraph{Is a retained divergence an error of the translation? (construct validity)} The reading of RQ1 as evidence against acceptances assumes that the C program is a valid reference on the inputs that are retained. The sanitizer check removes \rqUbInputsLibAllPct\% of the distinct divergent library inputs and \rqUbInputsExeAllPct\% of the executable ones, which shows how often generated inputs leave the defined behavior of the C program. On the inputs that are retained, the reference is valid only as far as the two sanitizers can tell, and \rqBThreeDYesExeAllNum{} of the \rqFPExeAllNum{} executable translations in $D$ lie in cases labeled as designed around undefined, implementation-defined or unspecified behavior (Section~\ref{sec:r-rq1}). A retained divergence in such a case can come from behavior that the sanitizers do not report or that the C standard leaves to the implementation. Generated library states are also not restricted to the preconditions of the function. When the C function rejects such a state through its own assertion and the translation does not, the inspection did not count the divergence as false evidence, since both programs were observed correctly. The case \code{027\_ctype\_ascii} shows a second form of the risk: on the same input, the runner and the harness reach different verdicts because they compare the output differently. The study limits this risk in two ways. A divergence is retained only when it is stable and the sanitizer check completes without a report, and an inspection of sampled translations in $D$ found no false evidence among \rqRNineFalseExeDen{} executable and \rqRNineFalseLibDen{} library translations that could be decided. In this inspection, each item was labeled in three passes of a language model, and the author resolved the disagreements and checked a subset of the agreed labels by hand; one retained divergence per translation was inspected. What remains is that a retained divergence is an observed difference under the comparison of this study. A defect of the C program, behavior that depends on the environment, or undefined behavior that the sanitizers miss can still produce one, so $D$ does not count translations that are proven to be semantically wrong.

\paragraph{Does the harness create or miss divergences? (internal validity)} All three results pass through one instrument, so a defect in it would bias them in an unknown direction, and a limited ability to detect would lower $D$. Table~\ref{tab:d-validation} relates each validation check to the risk it addresses and to the limitation that remains. Taken together, the checks give no sign that the harness produces retained divergences of its own: the rule tests pass, the unstable inputs of the C program compared with itself are all traced to the C program or to the time limit, a second compilation in place of the translation yields no retained divergence, and the inspected samples contain no false evidence. Detection is the weaker side. The harness finds every held-out translation again, but these translations diverge on a large share of their inputs, and it misses \sMutLibNeNotCaught{} of \sMutLibN{} non-equivalent library mutants. One gap lies in the comparison itself: for a library call that fails in the same way on both sides, the harness does not read the output streams, which concerns at most \rqLibGapUpperLibPct\% of the matching library inputs and can only hide divergences. The stability rule has a similar effect: a translation whose observation changes between runs counts as showing no retained divergence even when it always differs from C. \rqUnstableOnlyExeAllNum{} executable and \rqUnstableOnlyLibAllNum{} library translations have only unstable divergences, and the records do not show which side varied. The validation results indicate that the harness can miss behavioral divergences, particularly in libraries. These controls do not establish the extent of missed divergences in the study population or the overall direction of the measurement bias. Four checks departed from the plan. The pilot cases were not measured again on the final harness, and the held-out check was run after the main run on a machine of a different type. The check of the input domain has no record from before the main run and was run afterwards on the stored records; it found no retained divergence on an input outside the domain of the runner. The audit of the library divergences was not completed: its sample was drawn after the main run, and the manual reading of the sample was not carried out.

\begin{table}[!htb]
\centering
\small
\renewcommand{\arraystretch}{1.3}
\newcolumntype{L}[1]{>{\raggedright\arraybackslash}p{#1}}
\caption{Validation checks of the harness: the risk each check addresses, its result, and the limitation that remains.}
\label{tab:d-validation}
\begin{tabularx}{\linewidth}{@{}L{0.19\linewidth}L{0.23\linewidth}L{0.21\linewidth}>{\raggedright\arraybackslash}X@{}}
\toprule
Check & Risk addressed & Result & Remaining limitation \\
\midrule
Rule tests & Comparison or evidence rules implemented wrongly & \sUnitPass/\sUnitRun{} passed, \sUnitSkip{} skipped & Covers only the fixed cases of the tests \\
Replay of original vectors & Fault in build, execution or comparison & Agreement in all cases but one & In that case the runner and the harness compare the output differently, so the two checks can differ on one input \\
C against itself & Nondeterminism counted as divergence & \sVSixUnstable{} unstable inputs, \sVSixHarness{} caused by the harness & Checked on \sVSixInputsPerCase{} inputs per case \\
Second compilation of C as translation & Divergence produced by the measurement itself & \sCcExeRetained{} retained divergences & Run on the pilot cases only \\
Held-out failing vectors & Known divergences not found by generated inputs & \frc{rqHeldOutAll} targets rediscovered & Targets had frequent divergences; rare divergences remain untested \\
Non-equivalent mutants & Divergences missed by the search & \sMutExeK/\sMutExeN{} executable, \sMutLibK/\sMutLibN{} library & From \sMutExeNCases{} executable and \sMutLibNCases{} library pilot cases, labeled with model assistance; \sMutLibNeNotCaught{} library mutants missed \\
Inspection of samples (model-labeled, author-adjudicated) & False evidence in $D$, wrong outcome in $R_{\mathit{test}}$ & \rqRNineFalseExeNum{} false, \rqRTwoCheckWrongExeNum{} wrong among decided items & Samples only; \rqRNineLibUnclear{} library items could not be decided \\
\bottomrule
\end{tabularx}
\end{table}

\paragraph{Do the completed decisions stand for all decisions? (internal validity)} The proportions of RQ1 and RQ2 describe the decisions that reached the required outcome of the re-evaluation. Extending them to the full original populations would require assumptions about the decisions without a result, and the observed attrition gives reasons to question such assumptions. The attrition is not random. The main run was ended at a budget cut-off after the heaviest cases had filled the disks, so missing seeds and missing records account for \rqUnfinishedExeAllNum{} of the \rqLpExeAllNum{} executable acceptances without a result, and the \rqFFNotReexExeAllNum{} executable rejections without an outcome share one system and one failure. The study answers in two ways. The interval of RQ3 keeps every open decision in the reported rate, and the proportion of RQ1 can be bounded over all acceptances: between \rqFPBoundLoTightExeAllPct\% and \rqFPBoundHiExeAllPct\% for executables and between \rqFPBoundLoTightLibAllPct\% and \rqFPBoundHiLibAllPct\% for libraries, depending on whether only the acceptances with a recorded retained divergence or all acceptances without a result diverge. These bounds show the size of what is open. The causes of the attrition do not settle its direction. For \code{004\_nineality\_sieve} and \code{018\_stack\_buffer\_overflow\_loop1}, the operation log of the main run states that their divergence logs filled the disks, but many divergences need not yield retained divergences, since they can be unstable or set aside by the sanitizer check. For libraries, \rqLeaveIncompleteLibAllNum{} translations are incomplete because some inputs of their case could not be observed on any seed or machine, even after a retry, and the cause was not determined.

\paragraph{Do the results hold beyond this corpus? (external validity)} Any general statement about test-based evaluation of C-to-Rust translation assumes that the corpus and the systems are typical. The study covers one corpus, ten systems and the translations and decisions recorded at one commit of HARVEST, which matches the published paper; other branches of that repository record different decisions~\cite{harvest-results}. The cases are also not independent: one bundle contributes \nExePzeroP{} of the \nExePrimary{} executable and \nLibPzeroP{} of the \nLibPrimary{} library cases, and the five cases with the most translations in $D$ hold \rqAxTwoTopFiveLibPct\% of $D$ for libraries. The study reports the distribution over cases and draws no ranking of systems, since the systems were not evaluated on equal sets of translations. Counting by case--system pair does not remove this dependence. The results describe the published decisions on this corpus, and the thesis has no basis for a proportion of acceptances with a retained divergence in other corpora or for other systems.

\paragraph{Can the numbers be relied on and reproduced? (reliability)} A reader must be able to trace every number to its source and to obtain it again. The evidence rules, the metrics of RQ1 and RQ2 and the lower end of the interval of RQ3 were fixed before the main run, and every reported number is computed by scripts from the stored records, which are kept with the versions of the tools and of the harness. The decisions were reproduced with rustc~\sEnvRustcCorpus{}, as the corpus pins it, and the main run used rustc~\sEnvRustc{}; the replay of the original vectors in the main run agreed with the acceptances in all cases but one, where the cause lies in the comparison, so the change of compiler altered no verdict of an accepted translation on its original vectors. The part that depends on judgment is the inspection of samples. Its labels were assigned in three passes of a language model and adjudicated by the author, and passes of one model that rely on the same incomplete information can share an error.

\section{Implications for Evaluation Practice}\label{sec:d-implications}

The findings that remain after the threats of Section~\ref{sec:d-threats} concern how the evidence for a decision is gathered and how it is reported. Three implications follow for evaluations of C-to-Rust translation that rely on tests. Each is tied to one result of Chapter~\ref{ch:results}, and each inherits the limits of that result.

\paragraph{Check accepted translations on inputs beyond the test suite.} Within $P_{\mathit{test}}$, the acceptances of the original evaluation were consistent with the vectors on which they were decided, and the evidence against them came almost entirely from generated inputs: all translations in $D$ have a retained divergence on such an input. An evaluation that reports a pass rate can therefore strengthen it by running the accepted translations against the C program on additional inputs, with a check of the C side for undefined behavior. The result supports adding such a step. It does not prescribe the harness, the generators or the budget of this thesis, which were fixed for one corpus. The step also has a limited reach: under the controls of this study it missed some library mutants, so a translation that passes it has been searched and has not been shown to be equivalent.

\paragraph{Report the full denominator, the attrition and the resulting interval.} The success rate after reclassification could be stated only as an interval, \widthE{} percentage points wide for executables and \rqSRWidthTightLibAllPct{} for libraries, and for libraries the open decisions outweigh the smallest decrease. The size of a reported rate thus depends on decisions that were never examined, and a reader can judge this only when they are counted. A report of evaluation results should give the population from which the rate is computed, the pairs without a translation, the rejections at the build stage, the number of decisions that a later check actually reached, and the causes for which the others were not reached. Where decisions remain without a result, the rate is better stated as an interval. The interval of this thesis covers the open decisions only. Divergences that the inputs did not reach and undefined behavior that the sanitizers did not report lie outside it, so an interval of this kind marks the least uncertainty that a rate carries.

\paragraph{State the comparison rule and the conditions for using the reference.} A verdict depends on what is compared. When standard error was left out of the comparison, the unsupported executable rejections rose from \frc{rqFFExeAll} to \frc{rqRobFFNoStderrExeAll}, the one unsupported library rejection changed with the rule for floating-point fields, and in the case \code{027\_ctype\_ascii} the runner and the harness reached different verdicts on the same inputs. An evaluation should therefore publish which parts of the behavior it compares, which normalization and tolerance it applies, and how it treats executions on which the reference program has undefined behavior. In this study the rule moved only a few decisions, and the study cannot say how large the effect is on other corpora. The implication concerns the possibility of reproducing and comparing verdicts, whatever the size of the effect.

These implications are stated for the evaluation of translations. The results do carry information about the behavior of the translations of each system on this corpus, and Chapter~\ref{ch:results} reports it by system. They do not support an ordering of the systems or a general statement about the semantic quality of a translator, because the systems were re-evaluated on different sets of translations and the study is not a controlled comparison between them.
